\documentclass[]{article}
\usepackage[margin=1in]{geometry}
\usepackage{amsmath}
\usepackage{amssymb}
\usepackage{bbold}
\newcommand{\ket}[1]{\left\vert#1\right\rangle}
\newcommand{\bra}[1]{\left\langle#1\right\vert}
\newcommand{\norm}[1]{\|#1\|}
\newcommand{\alg}{\mathcal{A}}
\newcommand{\braket}[2]{\left\langle#1\vert#2\right\rangle}
\newcommand{\modulus}[1]{\left\vert#1\right\vert}
\newcommand{\ketbra}[2]{\ket{#1}\bra{#2}}
\newcommand{\set}[1]{\left\lbrace#1\right\rbrace}
\usepackage{fancyhdr}
\newcommand{\brak}[1]{\left\langle#1\right\rangle}
\newcommand{\com}[2]{\left[#1,#2\right]}
\newcommand{\acom}[2]{\left\lbrace#1,#2\right\rbrace}
\newcommand{\inner}[2]{\left\langle#1,#2\right\rangle}
\newcommand{\bd}[1]{\textbf{#1}}
\usepackage{color}

\begin{document}
\noindent \large Jeffrey M. Epstein\vspace{5pt}\hfill\small \today\\
National Institute of Standards and Technology \hfill jeffrey.m.epstein@gmail.com\\
Joint Center for Quantum Information and Computer Science \hfill jeffreymepstein.com
\\
\\
I am an NRC postdoctoral scholar at NIST Gaithersburg and the University of Maryland. I work on topics broadly related to quantum information, focusing in particular on quantum thermodynamics and on studies of the feasibility of near-term, low-energy tests of quantum gravity via quantum information-theoretic approaches.


\subsection*{Education}
\textbf{PhD}, Physics, December 2020\\
Dissertation: \textit{Statistical Mechanics of Transport Processes in Active Matter}\\
\textbf{MA}, Physics, December 2016\\
\textbf{University of California, Berkeley}; Berkeley, CA\\
\\
\textbf{MSc}, Perimeter Scholars International (PSI), June 2014\\
\textbf{Perimeter Institute for Theoretical Physics}; Waterloo, ON\\
\\
\textbf{AB}, Chemistry and Physics, May 2012\\
\textbf{Harvard College}; Cambridge, MA\\
\textit{magna cum laude} with high honors in field\\
secondary field, Mathematics; language citation, Chinese


\subsection*{Awards}
\begin{itemize}
	\item NIST NRC Postdoctoral Research Associateship, 2020
	\item National Defense Science and Engineering Graduate (NDSEG) Fellowship, 2016
	\item Herchel Smith-Harvard Undergraduate Science Research Program, 2010
\end{itemize}


\subsection*{Research}
I have worked in two main fields: quantum information and nonequilibrium transport.  In quantum information, I have established speed limits on information processing tasks derived from an extension of the Lieb-Robinson bound to time-dependent local Hamiltonians. More recently, I have demonstrated a connection between a class of nonlinear quantum amplifiers and the von Neumann model of measurement and worked on continuous-time quantum error correction. In 2014, as part of the Perimeter Scholars International masters program, I worked with Daniel Gottesman on magic state distillation. From August 2012-July 2013 I worked in the Quantum Information Group at IBM T.J. Watson research center, where I numerically studied the performance of randomized benchmarking, a method for assessing fidelities of quantum gates.\\
\\
Within nonequilibrium transport, I have focused on the emerging field of active matter. In particular, I have shown how active forces appear in the continuum equations of motion of a system of Active Brownian Particles, and I have derived Green-Kubo equations for generic two-dimensional non-equilibrium fluids, including an expression for the odd viscosity.

\subsection*{Scientific Publications}
\begin{enumerate}
\item \textit{Odd Diffusivity of Chiral Random Motion}. C Hargus, \textbf{JME}, KK Mandadapu. arXiv:2103.09958 (2021).

\item\textit{Quantum noise limits for a class of nonlinear amplifiers}. \textbf{JME}, KB Whaley, J Combes. Phys. Rev. A 103 (5), 052415 (2021).

\item \textit{Time reversal symmetry breaking and odd viscosity in active fluids: Green-Kubo and NEMD results}. C Hargus, K Klymko, \textbf{JME}, KK Mandadapu.  J. Chem. Phys. 152, 201102 (2020).

\item \textit{Time reversal symmetry breaking in two-dimensional non-equilibrium viscous fluids}. \textbf{JME}, KK Mandadapu. Phys. Rev. E 101, 052614 (2020).

\item \textit{Continuous quantum error correction for evolution under time-dependent Hamiltonians}. J Atalaya, S Zhang, MY Niu, A Babakhani, HCH Chan, \textbf{JME}, KB Whaley. 	arXiv:2003.11248 (2020).

\item  \textit{Statistical Mechanics of Transport Processes in Active Fluids II: Equations of Hydrodynamics for Active Brownian Particles}. \textbf{JME}, K Klymko, KK Mandadapu. J. Chem. Phys. 150, 164111 (2019).

\item \textit{Postponing the orthogonality catastrophe: efficient state preparation for electronic structure simulations on quantum devices}. NM Tubman, C Mejuto-Zaera, \textbf{JME}, D Hait, DS Levine, W Huggins, Z Jiang, JR McClean, R Babbush, M Head-Gordon, KB Whaley.	arXiv:1809.05523 (2018).



\item \textit{Quantum Speed Limits for Quantum Information Processing Tasks}. \textbf{JME}, KB Whaley. Phys. Rev. A 95, 042314 (2017).

\item \textit{Investigating the Limits of Randomized Benchmarking Protocols}. \textbf{JME}, AW Cross, E Magesan, and JM Gambetta. 
Phys. Rev. A 89, 062321 (2014) 

\item \textit{CD36 in the periphery and brain
	synergize in stroke injury in hyperlipidemia}. E Kim, M Febbraio, Y Bao, AT Tolhurst, \textbf{JME}, S Cho.  Annals of Neurology. 71(6) (2012)
\end{enumerate}







\subsection*{Teaching}
\begin{itemize}
	\item Graduate Student Instructor for Physics 112 (Introduction to statistical and thermal physics) UC Berkeley, Fall 2015
	
	\item Graduate Student Instructor for Physics 7b (Introductory thermodynamics and electromagnetism for scientists and engineers)  UC Berkeley, Fall 2014, Spring 2015
\end{itemize}


\subsection*{Conferences and Seminars}
\begin{enumerate}	
\item \textit{Nonlinear amplifiers for measurement of normal operators}. Talk presented to IBM Research - Almaden (2021).

\item \textit{Odd Transport in Active Systems}. Informal Statistical Physics Seminar, University of Maryland (2021).
	
\item \textit{Quantum Foundations Seminar}. Seminar series organized and taught at Berkeley (2020).

\item\textit{Active matter, time reversal symmetry breaking, and Onsager reciprocal relations}. Poster presented at Berkeley Statistical Mechanics Meeting (2020).
	
\item\textit{Rheology of 2D Active Fluids}. Berkeley Soft Matter Seminar (2019).	

\item \textit{Quantum noise limits for a class of nonlinear amplifiers}. Talk presented at APS March Meeting (2019).

\item \textit{Transport Processes in Active Fluids}. Talk presented at APS March Meeting (2019).



\item \textit{Continuum Mechanics of Active Brownian Particles}. Poster presented at Gordon Research Conference and Seminar: Complex Active and Adaptive Material Systems (2019).	

\item \textit{Speed Limits for Quantum Control of Local Spin Systems}. Talk presented at APS March Meeting (2018).


\item \textit{Speed Limits for Quantum Control of Local Spin Systems}. Talk presented at SQuInT (2017).


\item Lectures on Statistical Field Theories. School attended, Galileo Galilei Institute for Theoretical Physics, (2017).


\item \textit{Combinatorial Results on the Stabilizer Polytope}. Poster presented at SQuInT (2015).


\end{enumerate}


\end{document}
